% Source: Garre-Rubio--Lootens--Molnar, arXiv:2203.12563v3,
% REsubmission.tex lines 321--391, 431--490, and Appendix A, 2305 onwards.
% The reverse existence and uniqueness of biorthogonal decompositions remains
% separate from the sufficient local criteria below.

\section{Arbitrary boundary transport}
\label{sec:mpo_boundary_transport}

For $N\geq1$, an operator tensor $T$ and a state tensor $A$ define
\begin{align}
    O^N_{T,X}(\sigma,\tau)
     & =\tr\!\left(XT^{\sigma_1\tau_1}\cdots T^{\sigma_N\tau_N}\right),
    \label{eq:glm_boundary_operator}                                    \\
    \psi^N_{A,Y}(\sigma)
     & =\tr\!\left(YA^{\sigma_1}\cdots A^{\sigma_N}\right).
    \label{eq:glm_boundary_state}
\end{align}
No commutation condition is imposed on $X$ or $Y$. Cyclicity of trace
identifies $\psi^N_{A,Y}$ with the usual boundary parametrization
$\Gamma_N(Y)$ of the MPS ground space.

\begin{theorem}[Stacking arbitrary boundaries]
    \label{thm:glm_boundary_stack}
    \lean{MPOTensor.productBoundary,
        MPSTensor.mpvWithBoundary_eq_groundSpaceMap,
        MPOTensor.mpoWithBoundary_eq_physCloseN,
        MPOTensor.mpoWithBoundary_mulTensor,
        MPOTensor.mpvWithBoundary_actTensor,
        MPOTensor.mpoWithBoundary_mulVec_mpvWithBoundary}
    \leanok
    Write $R^{ik}=\sum_jP^{ij}\otimes Q^{jk}$ and
    $(T\cdot A)^i=\sum_jT^{ij}\otimes A^j$.
    For all boundary matrices and all lengths, including zero,
    \begin{align}
        O^N_{P,X}O^N_{Q,Y}    & =O^N_{R,X\otimes Y},
        \label{eq:glm_boundary_stack}                          \\
        O^N_{T,X}\psi^N_{A,Y} & =\psi^N_{T\cdot A,X\otimes Y}.
        \label{eq:glm_boundary_action}
    \end{align}
\end{theorem}
\begin{proof}
    \leanok
    Expand the product letters over the intermediate physical configuration.
    For the resulting word matrices $U,V$, use
    $(X\otimes Y)(U\otimes V)=XU\otimes YV$ and
    $\tr(XU\otimes YV)=\tr(XU)\tr(YV)$.
\end{proof}

For the diagram, write $R=PQ$ and $C=X\otimes Y$.
\begin{center}
    \tenkzkernel
    % Formula: (\ref{eq:glm_boundary_stack}).
    % Ink-to-index map: the first boxes carry X and Y and have no physical
    % legs; every other box is an operator letter. The two operator rows
    % contract the shared physical indices. Horizontal loops contract the
    % virtual indices through the displayed boundary matrices. Bonds are
    % authored explicitly: X and Y have no connecting physical edge.
    % Boundary signature on either side: (virtual-west=0 | virtual-east=0 |
    % physical-up=2 | physical-down=2). Dots denote the omitted sites.
    \begin{tenkz}[rows={op,op}, cols=4, boundary=periodic, bonds=none]
        \tn[skin=box,ports={180:virtual,0:virtual}]{X} &
        \tn[skin=mpo,ports={90:physical:$i_1$,270:physical}]{P} &
        \tn[skin=dots,ports={180:virtual,0:virtual}]{} &
        \tn[skin=mpo,ports={90:physical:$i_N$,270:physical}]{P} \\
        \tn[skin=box,ports={180:virtual,0:virtual}]{Y} &
        \tn[skin=mpo,ports={90:physical,270:physical:$k_1$}]{Q} &
        \tn[skin=dots,ports={180:virtual,0:virtual}]{} &
        \tn[skin=mpo,ports={90:physical,270:physical:$k_N$}]{Q}
        \tnwire{(1,1)}{(1,2)} \tnwire{(1,2)}{(1,3)} \tnwire{(1,3)}{(1,4)}
        \tnwire{(2,1)}{(2,2)} \tnwire{(2,2)}{(2,3)} \tnwire{(2,3)}{(2,4)}
        \tnwire{(1,2)}{(2,2)} \tnwire{(1,4)}{(2,4)}
    \end{tenkz}
    \quad $=$ \quad
    \begin{tenkz}[rows={op}, cols=4, boundary=periodic]
        \tn[skin=box,ports={180:virtual,0:virtual}]{C} &
        \tn[skin=mpo,ports={90:physical:$i_1$,270:physical:$k_1$}]{R} &
        \tn[skin=dots,ports={180:virtual,0:virtual}]{} &
        \tn[skin=mpo,ports={90:physical:$i_N$,270:physical:$k_N$}]{R}
    \end{tenkz}
\end{center}

\begin{definition}[Length-independent boundary closedness and compatibility]
    \label{def:glm_boundary_closedness}
    \lean{MPOTensor.IsBoundaryClosed, MPOTensor.IsBoundaryCompatible,
        MPOTensor.IsBoundaryClosed.mul_mem_range,
        MPOTensor.IsBoundaryCompatible.mulVec_mem_groundSpace}
    \leanok
    The closedness condition of \cite[Section~2]{GarreRubio2023MPOSym}
    requires, for every $X,Y\in\MN{\chi}$, a single $Z\in\MN{\chi}$
    such that $O^N_{T,X}O^N_{T,Y}=O^N_{T,Z}$ for every $N\geq1$.
    Compatibility requires, for every operator boundary $X\in\MN{\chi}$
    and state boundary $Y\in\MN{D}$, a single $Z\in\MN{D}$ such that
    $O^N_{T,X}\psi^N_{A,Y}=\psi^N_{A,Z}$ for every $N\geq1$.
    In both statements, the output boundary is chosen before the length.
    At each positive length, the operator range is therefore closed under
    multiplication, and every operator preserves the state ground space.
\end{definition}

\begin{theorem}[Separation by product boundaries]
    \label{thm:glm_boundary_separation}
    \lean{MPOTensor.productBoundary_single, MPOTensor.eq_iff_trace_productBoundary}
    \leanok
    Let $U,V$ be matrices on $\C^{D_1}\otimes\C^{D_2}$.
    They are equal if and only if
    $\tr((X\otimes Y)U)=\tr((X\otimes Y)V)$ for all matrices $X,Y$
    on the two factors.
\end{theorem}
\begin{proof}
    \leanok
    Taking $X=E_{ji}$ and $Y=E_{lk}$ identifies the
    $((i,k),(j,l))$ entries of $U$ and $V$, since
    $E_{ji}\otimes E_{lk}=E_{(j,l),(i,k)}$.
\end{proof}

\begin{theorem}[Linear boundary transport]
    \label{thm:glm_boundary_linear}
    \lean{MPOTensor.isBoundaryClosed_iff_exists_linearMap,
        MPOTensor.isBoundaryCompatible_iff_exists_linearMap}
    \leanok
    \uses{def:glm_boundary_closedness}
    Closedness is equivalent to the existence of a linear map
    $b:\MN{\chi^2}\to\MN{\chi}$ satisfying
    $O^N_{TT,C}=O^N_{T,b(C)}$ for every $C$ and $N\geq1$.
    Compatibility is equivalent to the existence of a linear map
    $b:\MN{\chi D}\to\MN{D}$ satisfying
    $\psi^N_{T\cdot A,C}=\psi^N_{A,b(C)}$ for every $C$ and $N\geq1$.
    These are the boundary maps at the start of
    \cite[Appendix~A]{GarreRubio2023MPOSym}.
\end{theorem}
\begin{proof}
    \leanok
    \uses{thm:glm_boundary_stack}
    Matrix units on a product bond space have the form $E_{ij}\otimes E_{kl}$.
    Choose an output boundary for each such matrix, simultaneously for all
    lengths, and extend linearly over this basis. Linearity of the
    boundary contraction gives the required identities. Conversely, apply
    the linear map to $X\otimes Y$ in
    (\ref{eq:glm_boundary_stack}) or (\ref{eq:glm_boundary_action}).
\end{proof}

\begin{definition}[Exact biorthogonal decomposition]
    \label{def:glm_boundary_decomposition}
    \lean{MPSTensor.IsBiorthogonalDecomposition}
    \leanok
    Let $B$ have bond dimension $D$, and let $A_c$ have bond dimension
    $D_c$ for each member of a finite index set. A biorthogonal decomposition
    consists of rectangular matrices $V_c:\C^D\to\C^{D_c}$ and
    $W_c:\C^{D_c}\to\C^D$ satisfying
    \begin{align}
        V_cW_c & =\Id_{D_c},         & V_cW_b & =0\quad(c\ne b),
        \label{eq:glm_boundary_biorthogonal}                     \\
        B^i    & =\sum_cW_cA_c^iV_c.
        \label{eq:glm_boundary_letter}
    \end{align}
    The index $c$ may contain a block label and its multiplicity index.
    These are the local fusion and action equations of
    \cite[Section~2]{GarreRubio2023MPOSym}.
\end{definition}

\begin{theorem}[Transport by an exact decomposition]
    \label{thm:glm_boundary_decomposition}
    \lean{MPSTensor.IsBiorthogonalDecomposition.sum_mul_sum,
        MPSTensor.IsBiorthogonalDecomposition.evalWord,
        MPSTensor.IsBiorthogonalDecomposition.trace_mul_evalWord,
        MPSTensor.IsBiorthogonalDecomposition.mpvWithBoundary,
        MPOTensor.mpoWithBoundary_mul_eq_sum_of_biorthogonalDecomposition,
        MPOTensor.mpoWithBoundary_mulVec_eq_sum_of_biorthogonalDecomposition}
    \leanok
    \uses{def:glm_boundary_decomposition}
    Under (\ref{eq:glm_boundary_biorthogonal}) and
    (\ref{eq:glm_boundary_letter}), each nonempty word $w$ satisfies
    \begin{align}
        B^w       & =\sum_cW_cA_c^wV_c,                      &
        \tr(XB^w) & =\sum_c\tr\!\left((V_cXW_c)A_c^w\right).
        \label{eq:glm_boundary_word}
    \end{align}
    In particular, an exact fusion decomposition of $PQ$ yields
    \begin{align}
        O^N_{P,X}O^N_{Q,Y}
         & =\sum_c O^N_{T_c,V_c(X\otimes Y)W_c},
        \label{eq:glm_boundary_fusion_transport}  \\
        O^N_{T,X}\psi^N_{A,Y}
         & =\sum_c\psi^N_{B_c,V_c(X\otimes Y)W_c}
        \label{eq:glm_boundary_action_transport}
    \end{align}
    for $N\geq1$, where the second identity uses an exact action
    decomposition of $T\cdot A$ into the tensors $B_c$.
\end{theorem}
\begin{proof}
    \leanok
    \uses{thm:glm_boundary_stack}
    Induct on the nonempty word. Multiplying two reconstructed products
    inserts $V_cW_b$; (\ref{eq:glm_boundary_biorthogonal}) removes all unequal
    intermediate labels. Cyclicity of trace moves $V_c$ to the left of
    the boundary. Apply (\ref{eq:glm_boundary_stack}) and
    (\ref{eq:glm_boundary_action}) to obtain the operator and state formulas.
\end{proof}

\begin{theorem}[Closedness from copies of one target]
    \label{thm:glm_boundary_closedness_of_decomposition}
    \lean{MPOTensor.isBoundaryClosed_of_biorthogonalDecomposition,
        MPOTensor.isBoundaryCompatible_of_biorthogonalDecomposition}
    \leanok
    \uses{def:glm_boundary_closedness,def:glm_boundary_decomposition}
    If $TT$ has an exact biorthogonal decomposition into copies of $T$,
    then $T$ is boundary closed. If $T\cdot A$ has such a decomposition
    into copies of $A$, then $T$ and $A$ are boundary compatible.
    In either case the output boundary is
    $Z=\sum_cV_c(X\otimes Y)W_c$.
\end{theorem}
\begin{proof}
    \leanok
    \uses{thm:glm_boundary_decomposition}
    Combine the sums in (\ref{eq:glm_boundary_fusion_transport}) or
    (\ref{eq:glm_boundary_action_transport}) by linearity in the boundary.
\end{proof}

\begin{theorem}[The empty-word condition]
    \label{thm:glm_boundary_empty_word}
    \lean{MPSTensor.IsBiorthogonalDecomposition.evalWord_of_complete}
    \leanok
    \uses{def:glm_boundary_decomposition}
    If in addition $\sum_cW_cV_c=\Id_D$, then
    $B^w=\sum_cW_cA_c^wV_c$ also holds for the empty word.
\end{theorem}
\begin{proof}
    \leanok
    \uses{thm:glm_boundary_decomposition}
    For $w=\varnothing$ the equality is the additional hypothesis.
    For nonempty words use (\ref{eq:glm_boundary_word}).
\end{proof}
